% !TEX program = pdflatex
\documentclass[11pt]{article}
\setlength{\textwidth}{6.5in}
\setlength{\oddsidemargin}{0pt}
\setlength{\evensidemargin}{0pt}
\setlength{\textheight}{9in}
\setlength{\topmargin}{0pt}
\addtolength{\topmargin}{-\headheight}
\addtolength{\topmargin}{-\headsep}
\ifdefined\pdfpagewidth \pdfpagewidth=8.5in \pdfpageheight=11in \fi
\ifdefined\pagewidth \pagewidth=8.5in \pageheight=11in \fi
\usepackage{amsmath,amssymb}
\usepackage{array}
\usepackage{booktabs}
\usepackage[tableposition=top]{caption}
\usepackage{longtable}
\usepackage{enumitem}
\usepackage{tikz}
\usepackage{pgfplots}
\pgfplotsset{compat=1.18}
\emergencystretch=3em

% ---- notation ---------------------------------------------------------------
\newcommand{\um}{\,$\mu$m}
\newcommand{\kms}{km\,s$^{-1}$}
\newcommand{\flux}{erg\,cm$^{-2}$\,s$^{-1}$}
\newcommand{\mA}{m\AA}
\newcommand{\sci}[2]{#1\times10^{#2}}
\newcommand{\ion}[2]{#1\,\textsc{#2}}
\newcommand{\nuc}[2]{$^{#1}$#2}
\newcommand{\lxl}{$L_X/L_\mathrm{bol}$}

% ---- front matter -----------------------------------------------------------
\title{Coronae, Pulses and Particles: Archival Tests for Artificial\\ Emission in Stellar Activity and High-Energy Surveys}
\author{Charles C. Norton}
\date{}

\begin{document}
\maketitle

\begin{abstract}
A technology on or near a star could add X-rays a corona cannot make, radio emission a
quiet star should not have, a pulse train, a flare clock, or particles from a direction
no astrophysical source occupies. We put each of these to the archive that measures it.
Of $137{,}896$ eROSITA coronal counterparts, none exceeds a tenth of its star's
luminosity in X-rays, and the $250$ above a hundredth are evolved stars, catalogued
cataclysmic, young or extragalactic sources and M dwarfs whose flares reach that level;
no slow rotator stands far above the rotation-activity relation, and among $63{,}043$
eROSITA point sources none holds its counts in a band where no continuum could put
them. No X-ray light curve of a star within $10$ to $20$\,pc holds a coherent period
beyond the noise. Radio emission from inactive stars, circular polarisation without a
counterpart and NVSS sources polarised beyond the synchrotron ceiling all trace to
background galaxies, active stars, pulsars or the diffuse emission of their fields, and
no GLEAM-X spectrum steps at a sub-band edge. No neutrino source sits at the $271$
systems within $10$\,pc or moves with Jupiter, Saturn, Mars or Venus, and ultra-high-energy
cosmic rays arrive in pairs and triplets as often as chance allows. In photon timing,
flares and pulses, every cluster of ultraviolet photons on a nearby star is a flare, no
flaring star keeps a clock better than a part in a thousand, planet hosts flare as their
matched stars do, and the dispersed pulses of the nearest stars' radio cadences appear
as often off the stars.
\end{abstract}

% =============================================================================
\section{Introduction}\label{sec:intro}
% =============================================================================

Stars are bright in every band, and a technology among them is most easily hidden
in their light. The places where it would show are the places where a star's own
physics sets a limit: a corona saturates at about a thousandth of the star's
luminosity~\cite{Pizzolato2003,Wright2011}; a quiet star's radio emission is faint and
unpolarised; synchrotron radiation polarises at most $72$ per cent; flares come at
random times; and neutrinos and the highest-energy cosmic rays come from cosmic
accelerators, not from nearby stars. A transmitter, a power plant or a beam would cross
one of these limits. Proposals have taken this direction for neutrinos~\cite{Learned2009}
and for structures around pulsars~\cite{Osmanov2016}. This paper tests the limits
against the archives that measure them, with a null drawn from the same data for each.

\subsection{Contributions}
\begin{enumerate}[leftmargin=*,itemsep=1pt]
\item Tests of stellar X-ray emission against the coronal ceiling, the rotation-activity
relation, the band structure of a continuum and coherent periodicity.
\item Tests of radio emission from inactive stars, of circular and linear polarisation
against what stars and synchrotron radiation allow, and of spectral edges and long-period
transients.
\item Tests of TeV neutrinos at nearby systems and moving with the planets, and of
ultra-high-energy cosmic rays arriving in pairs.
\item Tests of photon timing, flares and pulses: ultraviolet photon clusters, flare
clocks, superflare rates of planet hosts, box-shaped flares, dispersed pulses and pulsars
as hosts.
\end{enumerate}

% =============================================================================
\section{X-rays}\label{sec:xray}
% =============================================================================

\emph{The coronal ceiling.} A corona saturates near \lxl$=10^{-3}$, and flares raise it
for hours; a source holding a tenth of a star's luminosity in X-rays is not a corona.
The eRASS1 catalogue~\cite{Merloni2024} matched to Gaia by the HamStar
method~\cite{Schneider2022,Freund2024} gives $137{,}896$ coronal counterparts with a
distance and a colour, whose \lxl\ has percentiles $50$, $90$, $99$ and $99.9$ at
$-3.42$, $-2.85$, $-2.30$ and $-1.90$ dex. Two hundred and fifty exceed $10^{-2}$ and none
exceeds $10^{-1}$. Of the $250$, $153$ sit off Gaia's main sequence, $38$ carry a
cataclysmic, young, emission-line or extragalactic type on the counterpart or on a SIMBAD
object within $20''$, and $59$ are main-sequence M dwarfs of $3000$ to $3550$\,K at $50$
to $630$\,pc, whose flares reach $10^{-2}$ within eRASS1's day-long average. The one K
dwarf of the first pass, 1eRASS~J082623.3$-$703142 at $7.6\times10^{-2}$, is the
cataclysmic variable 1ES~0826$-$70.3, which SIMBAD files $11''$ from its Gaia counterpart.

\emph{Coronae without a dynamo.} A slowly rotating star has a weak corona, so X-rays far
above the rotation-activity relation of a slow rotator would need another source. For
$8287$ coronal counterparts of $3000$ to $6000$\,K with rotation periods from
TESS~\cite{Fetherolf2023}, Gaia DR3 and ZTF~\cite{Chen2020}, the Rossby number from the
turnover time of the colour's temperature~\cite{CranmerSaar2011} places each against the
relation of Wright et al.~\cite{Wright2011}: the residuals scatter by $0.82$\,dex below a
Rossby number of one, with $0.24$ per cent beyond $1.5$\,dex. Only $26$ stars reach a
Rossby number of one, since the period catalogues hold few slow rotators, and the one
beyond $1.5$\,dex is a ZTF RS~CVn binary at $1$\,kpc.

\emph{Line-only sources.} A continuum spreads its counts over the bands; an X-ray line
source at $0.5$ to $2$\,keV would hold them in one. Of the $63{,}043$ point sources of the
eRASS:3 catalogue with $80$ or more counts in four bands, the largest fraction in one
band has a median of $0.46$ and a $99$th percentile of $0.78$. Holding $85$ per cent in
one band at $3\sigma$ with every other below $5$ per cent are $13$ sources at
$0.2$--$0.5$\,keV, the supersoft class, $2$ at $2$--$5$\,keV, the absorbed class, and none
at $0.5$--$1$ or $1$--$2$\,keV, where no continuum could put them; the same test on the
released spectra is the step that remains.

\emph{Periodicities.} A transmitter pulsing in X-rays at seconds to minutes would show a
coherent period. In $209$ Chandra detections of the $50$ SIMBAD objects within $10$\,pc,
none reaches a false-alarm probability of $10^{-3}$ once the dither, the frame
sidebands and the red noise of flares are filtered; in $180$ XMM-Newton EPIC-pn
observations of stars within $20$\,pc, periods of $5$ to $600$\,s reach $0.01$ in four
against $3.6$ expected and $10^{-3}$ in none against $0.4$.

\emph{Nearby high-energy sources.} Of the $2626$ Gaia sources within $20$\,pc, carried to
eRASS1's epoch, the five above \lxl$=10^{-2.5}$ are M3 to M6 flare stars, and seven stars
lie inside the error ellipses of unassociated Fermi-LAT sources~\cite{Abdollahi2022}
against $7.9$ by chance.

% =============================================================================
\section{Radio}\label{sec:radio}
% =============================================================================

\emph{Inactive stars.} A quiet F to early K dwarf with no X-ray source and no activity
should be radio-silent. Against $605$ stellar radio emitters within $500$\,pc in
VLASS~\cite{Lacy2020}, $30$ inactive stars show $3$\,GHz emission, fewer than the $38$
chance alignments the shifted positions give; LoTSS-DR2~\cite{Shimwell2022} holds two
inactive single F to early K dwarfs and RACS-low~\cite{Hale2021} three, as many as chance
gives. The two brightest coincidences are background radio galaxies: $52.6$\,mJy at
$0.34''$ from Gaia DR3 2906691964510880128 by a steady steep spectrum from $0.86$ to
$3$\,GHz that no star gives, and $81$\,mJy at $1.2''$ from Gaia DR3 1283261877869284096,
resolved by LoTSS with a second WISE source beside the star.

\emph{Circular polarisation.} Coherent radio emission is circularly polarised, and in
stars it comes from activity. No RACS source of $10$ per cent circular polarisation at
signal-to-noise six lacks a counterpart among pulsars, active stars and published brown
dwarfs. Of the stars of V-LoTSS~\cite{Callingham2023} and the RACS
survey~\cite{Pritchard2021}, $23$ of the $31$ in eRASS1's half of the sky are detected in
X-rays at \lxl\ of $7.5\times10^{-5}$ to $3.6\times10^{-3}$, the rest are RS~CVn, BY~Dra
and flare stars by name and type, and none polarised above $30$ per cent has
\lxl\ below $10^{-5}$ and a period above $20$\,d.

\emph{The synchrotron ceiling.} Synchrotron radiation in a uniform field polarises at
most $72$ per cent, and a source beyond it needs a coherent emitter. Of the $26{,}973$
NVSS sources~\cite{Condon1998} above $10$\,mJy with polarised flux above $6$\,mJy, the
Rician-debiased fraction's $3\sigma$ lower bound has a median of $0.058$ and a $99$th
percentile of $0.51$; the $23$ above $0.72$ all have polarised fluxes at or above their
total fluxes, which no component can have, and lie in Fornax~A's lobes and on the Galactic
plane in Cygnus, Cassiopeia and Aquila, where the catalogue's polarised flux is the
field's diffuse emission. None sits at a pulsar.

\emph{Spectral edges.} A transmitter's band begins and ends abruptly. GLEAM-X~\cite{HurleyWalker2022}
gives $27{,}585$ sources with sixteen or more of its twenty $7.68$\,MHz sub-bands from
$72$ to $231$\,MHz; a step at a sub-band boundary fitted beside a curved power law gains a
$\chi^2$ whose median is $0.2$ and whose $99.9$th percentile is $71$, the flux-scale errors
common to all sources. The one source in the top thousandth with a step above $30$ per
cent, GLEAM-X~J082817.6$-$221400, shares the low-frequency beam with a $1.18$\,Jy source
$69''$ away, and the sub-band fits merge and split the pair as the beam shrinks.

\emph{Long-period transients and millimetre flares.} Of $16$ long-period radio
transients~\cite{Rea2026}, ten are given a natural nature by their papers, as white-dwarf
binaries, X-ray outbursts or rotators, and the other six hold no Gaia star at their
dispersion-measure distances beyond the rate at control positions. Of the $66$ stars of
the SPT-3G millimetre flares~\cite{Tandoi2024}, every one has an X-ray source except four
M5 to M8 emission-line dwarfs at $37$ to $49$\,pc, below eRASS1's limit.

% =============================================================================
\section{Particles}\label{sec:particles}
% =============================================================================

\emph{Neutrinos from nearby systems.} A power source or a beam of TeV neutrinos would
appear as a point source in IceCube's tracks. The unbinned point-source likelihood over
the ten-year sample~\cite{IceCube2021} at the $271$ systems within $10$\,pc finds none above
chance, and the stack of all has $p=0.26$; the median $90$ per cent limit on an $E^{-2}$
source is $1.6\times10^{29}$\,erg\,s$^{-1}$.

\emph{Neutrinos moving with the planets.} A source on a planet moves with it. The same
likelihood in each planet's frame, against $1000$ time scrambles that keep the events and
move the planet, gives test statistics of $0.9$ for Jupiter ($p=0.17$), $0.6$ for Saturn
($0.17$), $0.0$ for Mars ($0.42$) and $0.3$ for Venus ($0.23$), where the scrambles reach
$9$ to $13$; the $90$ per cent limits on $E^{-2}$ luminosity are $7\times10^{9}$\,W at
Venus, $10^{10}$\,W at Mars, $2\times10^{11}$\,W at Jupiter and $6\times10^{11}$\,W at
Saturn.

\emph{Cosmic-ray doublets.} Particles launched together arrive close in direction and
time. Among the $2635$ events above $32$\,EeV of the Pierre Auger Observatory's Phase
One~\cite{Auger2022}, pairs within $2^\circ$ number eight within ten days against a mean of
$5.7$ with the arrival times permuted ($p=0.21$) and two within a day against $0.56$
($p=0.11$), with no triplet within thirty days against $0.34$ expected; one pair's circle
holds Kapteyn's Star at $3.9$\,pc, its members $1.8$ and $2.1^\circ$ from it, beyond the
$0.9^\circ$ resolution. In the open data, the $2355$ inclined showers give seven pairs
against $5.1$ and the $22{,}751$ showers of the $1500$\,m array $569$ against $559$, with
$349$ triplets within thirty days against $361$.

% =============================================================================
\section{Photons in time}\label{sec:timing}
% =============================================================================

\emph{Ultraviolet photon clusters.} A pulsed ultraviolet source would gather photons in a
few milliseconds. Of $147$ GALEX near-ultraviolet eclipses of the $51$ stars within
$10$\,pc~\cite{Martin2005,Million2016}, nine hold a cluster at a false-alarm probability
below $0.01$ against $4.8$ from $1176$ control apertures; they are five events on Wolf~359,
Ross~128, Wolf~424, G~272$-$61 and GJ~725~A, all flare stars. GJ~725~A's is an isolated
$100$\,ms brightening of seven photons against $0.2$, clean of detector and field events and
fixed to the star in rebuilt detector coordinates; a census of $116$ eclipses finds $98$
such bursts across the detector, at $1.4\times10^{-6}$ per source-second, the mapped ones at
the detector's edge, and none on a catalogued star like GJ~725~A's. The far-ultraviolet
lists of $26$ stars hold no cluster, nor do $720$ control apertures, and HST's
photon-event lists of $108$ stars hold two below $10^{-3}$ against $3.7$ expected.

\emph{Flare clocks.} Flares come at random times, and a flare train at a fixed period
would be a clock. For $2304$ Kepler stars with ten or more catalogued
flares~\cite{YangLiu2019}, the H test~\cite{deJager2010} from $0.5$\,d to a third of the
span, against twenty draws of other stars' flare times in the same quarters, finds $22$
of the forty highest beating every draw away from rotation: one burst in three cycles,
two at the spacecraft's quarter or downlink spacing and $19$ locked in phase with
scatters of $0.004$ to $0.017$ of the period, none below the $10^{-3}$ of a clock. The two
long clocks on fast rotators are beats of two spot periods: KIC~7838383's $49.91$\,d from
$1.0472$ and $1.0685$\,d, and KIC~7746984's flare clock at half the $268$\,d beat of
$0.7398$ and $0.7377$\,d. Of $1426$ TESS stars~\cite{Yang2023}, nine of the forty highest beat
every draw, eight of them bursts at the $240$\,d spacing of TESS's pointings, and the
ninth, TIC~457100137 at $1.213$\,d, scatters by $0.0087$ of the period.

\emph{Superflares at planet hosts.} A civilization that dims or shields its star's
largest flares would leave its host short of superflares. Among $60{,}810$ TESS flares,
the $83$ flaring planet hosts have $28.3$ per cent of their flares above $10^{34}$\,erg
among those above $10^{33}$, against $35.0\pm6.0$ per cent for each host's ten nearest
non-hosts in temperature, radius, distance and sectors drawn $2000$ times ($p=0.14$), and
$48$ per cent of them carry a flare above $10^{34}$\,erg against $48\pm4$ per cent.

\emph{Box-shaped brightenings.} A flare rises fast and decays slowly, and a switched
source rises and falls alike. Among the $1806$ TESS $20$\,s light curves of stars within
$30$\,pc, $256$ brightenings at $7\sigma$ are fitted better by a box than by a flare,
against $263$ dips of the same shape: the rate of instrumental steps.

\emph{Dispersed pulses.} A pulse from a transmitter at a star is dispersed by the
interstellar medium in front of it and by nothing behind it. In Breakthrough Listen's
$349\,\mu$s products of $88$ scans of the nearest stars, broadband pulses at $7\sigma$ that
peak at a dispersion measure below $-1$, between $-1$ and $1$, and above $1$\,pc\,cm$^{-3}$
number $25$, $36$ and $9$ in the scans on the stars against $39$, $34$ and $12$ off them, so
the pulses are interference; folds at the Crab, Vela, B1937+21 and J0437$-$4715 periods
as each star sees them give $26$ of $176$ below $p=0.05$, but the off-source scans of a
control cadence fold as low, the red noise of the series.

\emph{Pulsars.} A structure around a pulsar would re-radiate its spin-down power in the
infrared. Of $1978$ pulsars with positions good to $1''$~\cite{Manchester2005}, $143$
have an AllWISE counterpart~\cite{Wright2010} against $98.6$ by chance, none of them ring-like
against $0.12$ per shifted copy, at a sensitivity of about $2$\,mJy at $12\,\mu$m for a
$300$\,K ring re-radiating $10^{32}$\,erg\,s$^{-1}$ at $1$\,kpc. The $45{,}116$ stars near the
antipodes of the hundred brightest pulsars, where a beacon timed to a pulsar would be
aimed, hold no result at $p<0.01$ in the coverage store.

\begin{table}[t]
\centering\small
\caption{The tests, their samples and results.}
\label{tab:tests}
\begin{tabular}{>{\raggedright\arraybackslash}p{0.22\textwidth} >{\raggedright\arraybackslash}p{0.3\textwidth} >{\raggedright\arraybackslash}p{0.4\textwidth}}
\toprule
Test & Sample & Result \\
\midrule
Coronal ceiling & $137{,}896$ eRASS1 stars & none above $10^{-1}$; the $250$ above $10^{-2}$ explained \\
Coronae without a dynamo & $8287$ stars with periods & one beyond $1.5$\,dex, an RS~CVn binary \\
Line-only X-ray sources & $63{,}043$ eRASS:3 sources & none at $0.5$--$2$\,keV \\
X-ray periodicities & $209$ Chandra, $180$ XMM detections & at the chance rate \\
Radio from inactive stars & VLASS, LoTSS, RACS & at the chance-alignment rate \\
Circular polarisation & RACS, V-LoTSS & active stars, pulsars, brown dwarfs \\
Synchrotron ceiling & $26{,}973$ NVSS sources & field emission in all $23$ \\
Spectral edges & $27{,}585$ GLEAM-X sources & one blend \\
TeV neutrinos & $271$ systems; four planets & at background \\
Cosmic-ray doublets & $2635$ events; open data & at the chance rate \\
Ultraviolet photon clusters & $51$, $26$ and $108$ stars & flares \\
Flare clocks & $2304$ Kepler, $1426$ TESS stars & beats, bursts and spacecraft periods \\
Superflares at hosts & $83$ hosts & as the matched stars \\
Dispersed pulses & $88$ scans & as often off the stars \\
Pulsars & $1978$ pulsars; $45{,}116$ stars & at the chance rate \\
\bottomrule
\end{tabular}
\end{table}

% =============================================================================
\section{Related work}\label{sec:related}
% =============================================================================

Learned, Pakvasa and Zee~\cite{Learned2009} proposed neutrinos for interstellar
signalling, and Osmanov~\cite{Osmanov2016} proposed infrared searches for structures around
pulsars. The tests here rest on the surveys' own work: the eROSITA catalogues and their
stellar identifications~\cite{Merloni2024,Schneider2022,Freund2024}, the rotation-activity
relation~\cite{Pizzolato2003,Wright2011}, the radio surveys and their stellar
samples~\cite{Lacy2020,Shimwell2022,Hale2021,Callingham2023,Pritchard2021,Condon1998,HurleyWalker2022},
the review of long-period transients~\cite{Rea2026}, the SPT-3G
flares~\cite{Tandoi2024}, IceCube's and the Pierre Auger Observatory's public
data~\cite{IceCube2021,Auger2022} and the flare catalogues of Kepler and
TESS~\cite{YangLiu2019,Yang2023}.

% =============================================================================
\section{Conclusion}\label{sec:conclusion}
% =============================================================================

No test finds emission a star or its surroundings cannot make. The X-ray sources above the
coronal ceiling are evolved, cataclysmic, young or flaring; the slow rotators sit on the
rotation-activity relation; no source holds its X-rays in a band no continuum could; no
nearby star pulses in X-rays. Radio emission from quiet stars arrives at the chance rate of
background sources, polarised emission belongs to active stars, pulsars or the diffuse
field, and no low-frequency spectrum steps. Neutrinos show no source at the nearest systems
or the planets, and cosmic rays pair as chance allows. Ultraviolet photon clusters are
flares, flare trains are beats and bursts, planet hosts flare as their twins do, box-shaped
flares are steps, and dispersed pulses are interference.

The reach of each test is its sample's. The X-ray tests cover eRASS1's half of the sky
at its depth and day-long average, so a source that flares beyond the ceiling for minutes
escapes it; the neutrino limits are $10^{22}$\,W at the nearest systems and $10^{10}$ to
$10^{12}$\,W at the planets; the polarisation and spectral tests reach the surveys'
flux limits; and the timing tests reach periods from milliseconds to the spans of the
photon lists and light curves.

\section*{Acknowledgements}

This work uses data from eROSITA, the soft X-ray instrument aboard SRG, through the eRASS1
and eRASS:3 releases; Gaia DR3 (ESA, Gaia DPAC); Chandra and XMM-Newton through their
archives; Fermi-LAT's 4FGL-DR3; VLASS (NRAO), LoTSS (LOFAR), RACS (ASKAP, CSIRO), NVSS
and GLEAM-X (MWA) through CIRADA, CASDA and VizieR; SPT-3G; the IceCube ten-year
point-source sample; the Pierre Auger Observatory's Phase One and open data; GALEX
through MAST and gPhoton and HST's photon-event lists; the Kepler and TESS flare
catalogues; Breakthrough Listen's open data; the ATNF Pulsar Catalogue and AllWISE; and
SIMBAD and VizieR at CDS, Strasbourg.

\begin{thebibliography}{99}\small\raggedright
\setlength{\itemsep}{1pt}

\bibitem{Pizzolato2003}
N.~Pizzolato, A.~Maggio, G.~Micela, S.~Sciortino and P.~Ventura, The stellar
activity-rotation relationship revisited: dependence of saturated and non-saturated X-ray
emission regimes on stellar mass for late-type dwarfs, \emph{Astron. Astrophys.}
\textbf{397}, 147 (2003).

\bibitem{Wright2011}
N.~J. Wright, J.~J. Drake, E.~E. Mamajek and G.~W. Henry, The stellar-activity-rotation
relationship and the evolution of stellar dynamos, \emph{Astrophys. J.} \textbf{743}, 48
(2011).

\bibitem{Learned2009}
J.~G. Learned, S.~Pakvasa and A.~Zee, Galactic neutrino communication, \emph{Phys. Lett.
B} \textbf{671}, 15 (2009).

\bibitem{Osmanov2016}
Z.~Osmanov, On the search for artificial Dyson-like structures around pulsars,
\emph{Int. J. Astrobiol.} \textbf{15}, 127 (2016).

\bibitem{Merloni2024}
A.~Merloni, G.~Lamer, T.~Liu, et al., The SRG/eROSITA all-sky survey: first X-ray
catalogues and data release of the western Galactic hemisphere, \emph{Astron.
Astrophys.} \textbf{682}, A34 (2024).

\bibitem{Schneider2022}
P.~C. Schneider, S.~Freund, S.~Czesla, et al., The eROSITA Final Equatorial-Depth Survey
(eFEDS): the stellar counterparts of eROSITA sources identified by machine learning and
Bayesian algorithms, \emph{Astron. Astrophys.} \textbf{661}, A6 (2022).

\bibitem{Freund2024}
S.~Freund, S.~Czesla, P.~Predehl, et al., The SRG/eROSITA all-sky survey: identifying the
coronal content with HamStar, \emph{Astron. Astrophys.} \textbf{684}, A121 (2024).

\bibitem{Fetherolf2023}
T.~Fetherolf, J.~Pepper, E.~Simpson, et al., Variability catalog of stars observed during
the TESS prime mission, \emph{Astrophys. J. Suppl.} \textbf{268}, 4 (2023).

\bibitem{Chen2020}
X.~Chen, S.~Wang, L.~Deng, R.~de~Grijs, M.~Yang and H.~Tian, The Zwicky Transient Facility
catalog of periodic variable stars, \emph{Astrophys. J. Suppl.} \textbf{249}, 18 (2020).

\bibitem{CranmerSaar2011}
S.~R. Cranmer and S.~H. Saar, Testing a predictive theoretical model for the mass loss
rates of cool stars, \emph{Astrophys. J.} \textbf{741}, 54 (2011).

\bibitem{Abdollahi2022}
S.~Abdollahi, F.~Acero, L.~Baldini, et al., Incremental Fermi Large Area Telescope fourth
source catalog, \emph{Astrophys. J. Suppl.} \textbf{260}, 53 (2022).

\bibitem{Lacy2020}
M.~Lacy, S.~A. Baum, C.~J. Chandler, et al., The Karl G. Jansky Very Large Array Sky
Survey (VLASS): science case and survey design, \emph{Publ. Astron. Soc. Pac.}
\textbf{132}, 035001 (2020).

\bibitem{Shimwell2022}
T.~W. Shimwell, M.~J. Hardcastle, C.~Tasse, et al., The LOFAR Two-metre Sky Survey. V.
Second data release, \emph{Astron. Astrophys.} \textbf{659}, A1 (2022).

\bibitem{Hale2021}
C.~L. Hale, D.~McConnell, A.~J.~M. Thomson, et al., The Rapid ASKAP Continuum Survey
paper II: first Stokes I source catalogue data release, \emph{Publ. Astron. Soc. Aust.}
\textbf{38}, e058 (2021).

\bibitem{Callingham2023}
J.~R. Callingham, T.~W. Shimwell, H.~K. Vedantham, et al., V-LoTSS: the circularly
polarised LOFAR Two-metre Sky Survey, \emph{Astron. Astrophys.} \textbf{670}, A124 (2023).

\bibitem{Pritchard2021}
J.~Pritchard, T.~Murphy, A.~Zic, et al., A circular polarization survey for radio stars
with the Australian SKA Pathfinder, \emph{Mon. Not. R. Astron. Soc.} \textbf{502}, 5438
(2021).

\bibitem{Condon1998}
J.~J. Condon, W.~D. Cotton, E.~W. Greisen, et al., The NRAO VLA Sky Survey, \emph{Astron.
J.} \textbf{115}, 1693 (1998).

\bibitem{HurleyWalker2022}
N.~Hurley-Walker, T.~J. Galvin, S.~W. Duchesne, et al., GaLactic and Extragalactic All-sky
Murchison Widefield Array survey eXtended (GLEAM-X) I: survey description and initial
data release, \emph{Publ. Astron. Soc. Aust.} \textbf{39}, e035 (2022).

\bibitem{Rea2026}
N.~Rea, N.~Hurley-Walker, M.~Caleb, et al., Long period transients (LPTs): a comprehensive
review, arXiv:2601.10393 (2026).

\bibitem{Tandoi2024}
C.~Tandoi, S.~Guns, A.~Foster, et al., Flaring stars in a non-targeted mm-wave survey with
SPT-3G, \emph{Astrophys. J.} \textbf{972}, 6 (2024).

\bibitem{IceCube2021}
IceCube Collaboration, R.~Abbasi, M.~Ackermann, et al., IceCube data for neutrino
point-source searches years 2008--2018, arXiv:2101.09836 (2021).

\bibitem{Auger2022}
Pierre Auger Collaboration, P.~Abreu, M.~Aglietta, et al., Arrival directions of cosmic
rays above 32 EeV from Phase One of the Pierre Auger Observatory, \emph{Astrophys. J.}
\textbf{935}, 170 (2022).

\bibitem{Martin2005}
D.~C. Martin, J.~Fanson, D.~Schiminovich, et al., The Galaxy Evolution Explorer: a space
ultraviolet survey mission, \emph{Astrophys. J.} \textbf{619}, L1 (2005).

\bibitem{Million2016}
C.~Million, S.~W. Fleming, B.~Shiao, et al., gPhoton: the GALEX photon data archive,
\emph{Astrophys. J.} \textbf{833}, 292 (2016).

\bibitem{YangLiu2019}
H.~Yang and J.~Liu, The flare catalog and the flare activity in the Kepler mission,
\emph{Astrophys. J. Suppl.} \textbf{241}, 29 (2019).

\bibitem{deJager2010}
O.~C. de Jager and I.~B\"usching, The H-test probability distribution revisited: improved
sensitivity, \emph{Astron. Astrophys.} \textbf{517}, L9 (2010).

\bibitem{Yang2023}
Z.~Yang, L.~Zhang, et al., Properties of flare events based on light curves from the TESS
survey, \emph{Astron. Astrophys.} \textbf{669}, A15 (2023).

\bibitem{Manchester2005}
R.~N. Manchester, G.~B. Hobbs, A.~Teoh and M.~Hobbs, The Australia Telescope National
Facility pulsar catalogue, \emph{Astron. J.} \textbf{129}, 1993 (2005).

\bibitem{Wright2010}
E.~L. Wright, P.~R.~M. Eisenhardt, A.~K. Mainzer, et al., The Wide-field Infrared Survey
Explorer (WISE): mission description and initial on-orbit performance, \emph{Astron. J.}
\textbf{140}, 1868 (2010).

\end{thebibliography}

\end{document}
